% Pursuing Stacks (cotes 134-2 to 134-8): archive cote 134-5, site page 3 (PDF page 4), typescript page 190.
% typescript. Draft compared: ps4.tex:98-173.
% AI-transcribed by Claude Opus 5.5 (claude-opus-5-5, adaptive thinking, effort high), Message Batches API, batch ps2 (msgbatch_01NVwr1Yuo8qJf5rykeGBXPK), 2026-11-12; not yet checked by eye.
\scanpage{134-5/3}
\mspage{190}

\diagram{\begin{tikzcd}
A\hat{} \arrow[r, "i_A"] \arrow[d, swap, "\hand{\gamma_A}"] & (Cat) \arrow[d, "\hand{\gamma_{\underline{W}} \text{ or } \gamma}"] \\
(Hot_A) \arrow[r, "\overline{i}_A"] & (H\unsure{\hand{\hat{o}}}t(\underline{W}))
\end{tikzcd}}
(3) \hfill ,

we denote by
\[ \hand{\varphi}_A= \hand{\gamma i_A = \overline{i}_A\gamma_A} : A\hat{} \longrightarrow (Hot(\underline{W})) \]
the corresponding composition.

Coming back to the case of a functor i and corresponding $i^*$
((1) and (2)), the functors $i^*$ , $i_A$ , $i_B$ do \uline{not} give rise to a commu-%
tative triangle, but to a triangle \uline{with commutation morphism} $\hand{\lambda}_i$ :
\diagram{\begin{tikzcd}[column sep=small]
B\hat{} \arrow[rr, "i^*"] \arrow[dr, swap, "i_B", ""{name=iB, right}] & & A\hat{} \arrow[dl, "i_A", ""{name=iA, left}] \\
& (Cat) \arrow[bend right=10, from=iA, to=iB, swap, "\hand{\lambda_i}"] &
\end{tikzcd}}
\hfill ,

i.e. for any F in $B\hat{}$ , we get a map \xout{\illegible{1}}
\[ (4) \qquad \hand{\lambda}_i(F) : i_Ai^*(F) = A_{/i^{*}(F)} \longrightarrow i_B(F) = B_{/F} \quad , \]
functorial in F for variable F. \unsure{T\hand{h}he|The} first hand side of (4), also
written simply $A_{/F}$ when there is no ambiguity for i , \xout{ixxth} can be
interpreted as the category of pairs
\[ (a,p) \ , \quad a \in \mathrm{Ob}\, A \ , \quad p: i(a) \to F \quad , \]
where p is a map in $B\hat{}$, B identified as usual to a full subcategory of
$B\hat{}$ (hence i(a) identified with an \unsure{object|obbject} of $B\hat{}$ ). The \xout{functor} \inserted{map} $\hand{\lambda}_i(F)$
for fixed F is of course the \xout{functxxxxxx} functor
\[ (a,p) \longmapsto (i(a),p) \quad : \quad \hand{A_{/F} \longrightarrow B_{/F} .} \]

The topological significance of course is clear: interpreting i as
defining a"map" or morphism of the corresponding topoi $A\hat{}$ and $B\hat{}$, having
$i^*$ as inverse image functor, an object F of $B\hat{}$ gives rise to an \uline{induced}
\uline{topos} $B\hat{}_{/F} \hand{\overset{\approx}{\longrightarrow}} (B_{/F})\hat{}$ , and the restriction of the ``topos above'' $A\hat{}$
to the induced $B\hat{}_{/F}$ , or equivalently the result of base change
$B\hat{}_{/F} \to B\hat{}$ , gives rise to the \inserted{induced} morphism of topoi $A\hat{}_{/i^*(F)} \longrightarrow B\hat{}_{/F}$ ,
represented precisely by the map $\hand{\lambda}_i(F)$ in (Cat).

The condition \xout{taht} of $\underline{W}$-asphericity on i may be expressed in many-%
fold ways, as properties o\over{f}{f} either one of the three aspects
\[ \hand{\lambda}_i \ , \quad i^* \ , \quad i_Ai^* \]
of the situation created by i , \deleted{and} \inserted{with respect to} the localizing sets $\underline{W}$ , $\underline{W}_A$ , $\underline{W}_B$ \hand{,} \inserted{or \unsure{to} the notion of aspheric objects.} \margin{$\checkmark$}
As $\underline{W}_A$ is defined in terms of $\underline{W}$ as just the inverse image of the latter
by $i_{\over{A}{A}}$ , \inserted{and the same for aspheric objects,} \hand{(}it turns out that each of the conditions we \xout{ax}
\margin{3}
